
\def\PUDcodigo{ECA208}

\if\COMPILAR0
    \def\PUDcodacad{04507.10}
\else
    \def\PUDcodacad{04506.11}
\fi

\def\PUDdisciplina{Lógica de Programação}

%NOVOS ---------------------------------------------------
\def\PUDhorasTeoria{40}
\def\PUDhorasPratica{40}

\if\COMPILAR0
   \def\PUDinterECA{ \hyperref[PUD:ECA209]{Linguagem de Programação (04507.14)} }
\else
   \def\PUDinterECA{ \hyperref[PUD:ECA209]{Linguagem de Programação (04506.15)} }
\fi

\def\PUDatuaisECA{
- Aplicações com tecnologias atuais
}

\def\PUDrecursos{Projetor multimídia; Quadro branco e pincel; Aulas em laboratório de informática.}

%----------------------------------------------------------

\def\PUDsemestre{S2}

\def\PUDhoras{80}

\def\PUDcreditos{4}

\def\PUDprerequisito{---}

\def\PUDpreacad{---}

\def\PUDobjetivos{Compreender os fundamentos de lógica de programação e desenvolvimento de programas estruturados.  Ler e escrever programas de computador utilizando linguagem de algoritmos e linguagem de alto nível. Além de verificar o contexto de aplicações regionais e interdiciplinariedade com diversas disciplinas presentes na matriz do curso.} %A ÚLTIMA FRASE É NOVA.

\def\PUDementa{Introdução ao conceito de algoritmo. Desenvolvimento de algoritmos. Os conceitos de variáveis, tipos de dados, constantes, operadores aritméticos, expressões, atribuição, estruturas de controle (atribui- ção, seqüência, seleção, repetição). Vetores e Matrizes. Procedimentos, funções e bibliotecas. Metodologias de desenvolvimento de programas. Representação gráfica e textual de algoritmos. Estrutura e funcionalidades básicas de uma linguagem de programação procedimental. Implementação de algoritmos através da linguagem de programação introduzida. Depuração de Código e Ferramentas de Depuração.
%Introdução ao conceito de algoritmo.  Desenvolvimento de algoritmos.  Os conceitos de variáveis, tipos de dados, constantes, operadores aritméticos, expressões, atribuição, estruturas de controle (atribuição, seqüência, seleção, repetição).  Metodologias de desenvolvimento de programas.  Representação gráfica e textual de algoritmos.  Estrutura e funcionalidades básicas de uma linguagem de programação procedimental.  Implementação de algoritmos através da linguagem de programação introduzida.  Depuração de Código e Ferramentas de Depuração, Módulos (Procedimentos, Funções, Unidades ou Pacotes, Bibliotecas), Recursividade, Ponteiros e Alocação Dinâmica de Memória, Estruturas de Dados Heterogêneas (Registros ou Uniões, Arrays de Registros), Arquivos: Rotinas para manipulação de arquivos, Arquivos texto, Arquivos Binários, Arquivos de Registros.
}

\def\PUDconteudo{ & Técnicas de Elaboração de Algoritmos e Fluxogramas;  Algoritmos;   Fluxograma. 
\\ 
 & Linguagem C;     Constantes: numérica, lógica e literal;             Variáveis: formação de identificadores, declaração de variáveis, comentários e comandos de atribuição;     Expressões e operadores aritméticos, lógicos, relacionais e literais, prioridade das operações;     Comandos de entrada e saída;     Estrutura seqüencial, condicional e de repetição. 
\\ 
 & Estrutura de dados;     Variáveis compostas homogêneas unidimensionais (vetores).     Variáveis compostas homogêneas multidimensionais (matrizes).     %Variáveis compostas heterogêneas (registros).     Arquivos. 
\\ 
 & Modularização;     Procedimentos e funções;    
 %Passagens de parâmetros;     Regras de escopo. 
 Aplicações com tecnologias atuais. 
%\\ 
 %& Interfaces;     Porta paralela no PC;     Porta Serial RS232. 
 }

\def\PUDmetodo{Aulas expositórias que podem ser teóricas e/ou práticas, onde as práticas no laboratório serão marcadas no decorrer da disciplina.}

\def\PUDavalia{A avaliação será feita com aplicação de provas teóricas e/ou práticas, além da possibilidade de inclusão de trabalhos e seminários no decorrer da disciplina.}

\def\PUDbibbasica{ & \href{www.biblioteca.ifce.edu.br/index.asp?codigo_sophia=62925}{Ziviani}, N.  Projeto de algoritmos: com implementações em Java e C++.  São Paulo, SP : Cengage Learning, 2015.  621p.  ISBN: 9788522105250.
\\ 
 & \href{www.biblioteca.ifce.edu.br/index.asp?codigo_sophia=24429}{Beneduzzi}, Humberto Martins.  Lógica e Linguagem de Programação: Introdução ao Desenvolvimento de Software.  144p.  ISBN: 9788563687111. 
\\ 
 & \href{www.biblioteca.ifce.edu.br/index.asp?codigo_sophia=24502}{Souza}, M. A. F.; GOMES, M. M.; SOARES, M. V.; CONCILIO, R.  Algoritmos e Lógica de Programação.  São Paulo, SP : Cengage Learning, 2008.  214p.  ISBN: 8522104646. }

\def\PUDbibcomplementar{ & \href{www.biblioteca.ifce.edu.br/index.asp?codigo_sophia=24202}{Ziviani}, Nivio.  Projeto de algoritmos: com implementação em Pascal e C.  3º ed.  São Paulo, SP: Cengage Learning, 2011.  639p.  ISBN: 9788522110506.
\\ 
 & \href{www.biblioteca.ifce.edu.br/index.asp?codigo_sophia=24181}{Schildt}, H.  C: completo e total.  3º ed.  São Paulo, SP : Pearson Makron Books, 1997.  827p.  ISBN: 9788534605953.
\\ 
 &  \href{www.biblioteca.ifce.edu.br/index.asp?codigo_sophia=65446}{Cormen}, T. H.  Algoritmos: teoria e prática.  2º ed.  Rio de Janeiro, RJ: Elsevier, 2012.  926p.  ISBN: 9788535236996.
\\
 &  \href{www.biblioteca.ifce.edu.br/index.asp?codigo_sophia=56907}{Ascencio}, Ana Fernanda Gomes; Campos, Edilene Aparecida Veneruchi de.  Fundamentos da Programação de Computadores.  E-book.  Pearson.  588p.  ISBN: 9788564574168.
\\
 &  \href{www.biblioteca.ifce.edu.br/index.asp?codigo_sophia=40651}{Forbellone}, André Luiz Villar.  Lógica de programação: a construção de algoritmos e estruturas de dados.  3º ed.  São Paulo, SP: Pearson, 2013.  218p.  ISBN: 9788576050247. }

\def\PUDautor{Geraldo Luis Bezerra Ramalho}

\def\PUDdata{2013-04-17}

\def\PUDrevisao{3}

\def\PUDrevisor{Antonio Barbosa de Souza Júnior}

\def\PUDdatarevisao{2019-05-15}

\PUDcorpo
